% Corrected French for printed p.172; departures are append-only below.
\SourcePage{172}

\LayerTarget{tohoku:eq:direct-image-support}
\[
\operatorname{supp.}_{X}(\varphi)
=\text{adhérence de }f(\operatorname{supp.}_{Y}\varphi)
\]
pour toute $\varphi\in\Gamma(f_\Psi(F))=\Gamma_{\Psi(X)}(F)$ (dont on
considère les supports dans $X$ et dans $Y$, qu'on distingue dans la notation
en mettant l'espace en indice au symbole supp.).

\LayerTarget{tohoku:def:pushforward-adapted-support-antifilter}
Soit alors $\Phi$ un antifiltre de parties fermées dans $X$, désignons par
$\Psi'$ l'antifiltre des parties fermées $A$ de $Y$ qui sont dans $\Psi(X)$,
et telles que l'adhérence de $f(A)$ soit dans $\Phi$. Alors les formules
précédentes impliquent la formule :

\LayerTarget{tohoku:eq:adapted-support-sections}
\[
\Gamma_\Phi(f_\Psi(F))=\Gamma_{\Psi'}(F).
\]

\LayerTarget{tohoku:def:adapted-antifilters}
(Les familles $\Phi$ et $\Psi$ sont dites \textit{adaptées (relativement à
$f$)} si on a $\Psi=\Psi'$. C'est le cas par exemple si $\Phi$ et $\Psi$ sont
formés de toutes les parties fermées de $X$ resp. $Y$).

\LayerTarget{tohoku:functor:relative-direct-image-left-exact}
\textit{$f_\Psi$ est un foncteur exact à gauche} de la catégorie $C^Y$ des
faisceaux abéliens sur $Y$ dans la catégorie $C^X$, de plus la formule
précédente désigne un isomorphisme \textit{fonctoriel}. On peut l'écrire

\LayerTarget{tohoku:eq:adapted-support-functor-factorization}
\[
\Gamma_{\Psi'}=\Gamma_\Phi f_\Psi.
\]
Nous voulons lui appliquer le théorème 2.4.1, pour ceci il faut donner des
conditions moyennant lesquelles $f_\Psi$ transforme un faisceau injectif en
un faisceau $\Gamma_\Phi$-acyclique.

\LayerTarget{tohoku:lemma:3.7.1}
\noindent\textsc{Lemme 3.7.1.}
\begin{description}
\LayerTarget{tohoku:subitem:3.7.1:1}
\item[\textit{1)}] \textit{Si $\Psi$ est l'ensemble de toutes les parties
fermées de $Y$, $f_\Psi$ transforme faisceaux injectifs en faisceaux
injectifs.}
\LayerTarget{tohoku:subitem:3.7.1:2}
\item[\textit{2)}] \textit{Si $\Phi$ est paracompactifiante, alors $f_\Psi$
transforme faisceaux flasques en faisceaux $\Phi$-mous.}
\end{description}

\LayerTarget{tohoku:proof:3.7.1:1}
\noindent\textsc{Démonstration.} 1) Supposons que $F$ soit injectif ; soit,
pour tout $y\in Y$, $M(y)$ un groupe abélien injectif contenant $F(y)$, et
soit $M$ le faisceau produit défini par les $M(y)$ (cf. 3.1). $F$ est un
sous-faisceau de $M$, donc un facteur direct de $M$ puisqu'il est injectif,
donc $f_\Psi(F)$ est un facteur direct de $f_\Psi(M)$, et il suffit de prouver
que $f_\Psi(M)$ est injectif.

\LayerTarget{tohoku:def:product-direct-image-sheaf}
Or pour tout $x\in X$, soit $N(x)$ le groupe produit des $M(y)$ pour
$y\in f^{-1}(x)$, il résulte aussitôt des définitions que $f(M)$ est le
faisceau produit $N$ défini par la famille des $N(x)$. Or chaque $N(x)$ est un
groupe abélien injectif comme produit de groupes injectifs, donc $N$ est
injectif (prop. 3.1.2) donc $f(M)$ est injectif.

\LayerTarget{tohoku:proof:3.7.1:2:p1}
2) Supposons $F$ flasque, prouvons que $f_\Psi(F)$ est $\Phi$-mou. Soit donc
$g$ une section de $f_\Psi(F)$ sur un $B\in\Phi$, on cherche une
$h\in\Gamma_\Phi(f_\Psi(F))=\Gamma_{\Psi'}(F)$ dont la restriction à $B$
soit $g$. Comme $\Phi$ est paracompactifiante, $B$ admet un voisinage
$B'\in\Phi$ paracompact, donc $g$ est la restriction d'une section $g'$ de
$f_\Psi(F)$ définie sur un voisinage convenable $U\subset B'$ de $B$. $B'$
étant normal il existe un voisinage \textit{fermé} $B_1$ de $B$ contenu dans
$U$, soit $U_1$ son intérieur.

\LayerTarget{tohoku:construction:flabby-extension-section}
Considérons $g'$ comme un élément de
$\Gamma_{\Psi(U)}(f^{-1}(U),F)$, son support $A$ est une partie fermée de
$f^{-1}(U)$, donc $A\cap f^{-1}(B_1)$ est une partie fermée de $Y$, donc son
complémentaire dans $Y$ est ouvert. L'intersection de cet ouvert avec l'ouvert
$f^{-1}(U_1)$ étant contenue dans $\mathrm C A$, $g'$ y est nulle, donc il
existe une section $g_1$ de $F$ sur l'ouvert réunion de $f^{-1}(U_1)$ et de
$\mathrm C(A\cap f^{-1}(B_1))$ qui coïncide avec $g'$ sur le premier et est
nulle sur le second. Enfin, $F$ étant flasque, il existe une section $h$ de
$F$ sur $Y$

\TohokuContinuation{173}
